% Biên dịch bằng LuaLaTeX (đã đặt sẵn trong latexmkrc)
\documentclass[11pt,a4paper]{article}

% ---------- Tiếng Việt + font ----------
\usepackage{fontspec}
\usepackage[vietnamese,provide=*]{babel}
\setmainfont{Noto Serif}
\setsansfont{Noto Sans}
\setmonofont{Noto Sans Mono}[Scale=0.9]

% ---------- Bố cục ----------
\usepackage[margin=2.5cm]{geometry}
\usepackage[parfill]{parskip}        % tách đoạn bằng một dòng trống, không thụt đầu dòng
\usepackage{titlesec}
\titleformat{\section}{\Large\bfseries\sffamily\color{accent}}{\thesection}{0.6em}{}
\titleformat{\subsection}{\large\bfseries\sffamily}{\thesubsection}{0.6em}{}
\titleformat{\subsubsection}{\normalsize\bfseries\sffamily}{\thesubsubsection}{0.6em}{}

% ---------- Thành phần thường dùng ----------
\usepackage{amsmath,amssymb}
\usepackage{graphicx}
\graphicspath{{images/}}             % ảnh để trong thư mục images/
\usepackage{float}                   % [H] = đặt ảnh/bảng đúng chỗ
\usepackage{booktabs,tabularx}
\usepackage{enumitem}
\setlist{nosep, leftmargin=*}
\newlist{checklist}{itemize}{2}
\setlist[checklist]{label=$\square$}
\usepackage{xcolor}
\definecolor{accent}{HTML}{1F4E79}
\usepackage[most]{tcolorbox}
\usepackage[numbers]{natbib}
\usepackage[colorlinks=true,linkcolor=accent,citecolor=accent,urlcolor=accent]{hyperref}

% ---------- Khung code ----------
\usepackage{listings}
\definecolor{codebg}{HTML}{F6F8FA}
\lstset{
  basicstyle=\ttfamily\small,
  backgroundcolor=\color{codebg},
  frame=single, rulecolor=\color{black!15},
  numbers=left, numberstyle=\tiny\color{black!45},
  keywordstyle=\color{accent}\bfseries,
  commentstyle=\color{black!55}\itshape,
  stringstyle=\color{red!60!black},
  breaklines=true, showstringspaces=false,
  tabsize=4, columns=fullflexible, keepspaces=true,
  xleftmargin=1.5em, framexleftmargin=1.5em,
}

% ---------- Hộp ghi chú ----------
\newtcolorbox{keyidea}[1][Ý chính]{colback=blue!4, colframe=accent, fonttitle=\bfseries\sffamily, title=#1}
\newtcolorbox{question}[1][Câu hỏi]{colback=orange!5, colframe=orange!70!black, fonttitle=\bfseries\sffamily, title=#1}
\newtcolorbox{warning}[1][Lưu ý]{colback=red!4, colframe=red!60!black, fonttitle=\bfseries\sffamily, title=#1}

\newcommand{\todo}[1]{\textcolor{red}{[TODO: #1]}}

% ---------- Tiêu đề ----------
\title{\sffamily\bfseries Ghi chú: Research Engineer}
\author{Lâm Quang Anh Quân}
\date{\today}

\begin{document}
\maketitle
\tableofcontents

% =====================================================================
\section{Cách đọc paper}

Đoạn văn đầu tiên. Để sang đoạn mới, chỉ cần để \textbf{một dòng trống}
giữa hai đoạn. Xuống dòng một lần trong file \texttt{.tex} vẫn tính là cùng một đoạn.

Đây là đoạn thứ hai. Chữ có thể \textbf{in đậm}, \emph{in nghiêng},
\underline{gạch chân} hoặc \texttt{chữ code}. Chú thích cuối trang đặt như sau\footnote{Ví dụ chú thích cuối trang.}.
Trích dẫn tài liệu: phương pháp đọc ba lượt~\citep{keshav2007read}.

\subsection{Phương pháp đọc ba lượt}

\subsubsection{Lượt 1: đọc lướt}

Gạch đầu dòng, có thể lồng nhiều cấp:
\begin{itemize}
  \item Đọc tiêu đề, abstract, phần giới thiệu.
  \item Xem các hình và bảng.
  \begin{itemize}
    \item Gạch đầu dòng cấp 2.
    \begin{itemize}
      \item Gạch đầu dòng cấp 3.
    \end{itemize}
  \end{itemize}
  \item Đọc phần kết luận.
\end{itemize}

\subsubsection{Lượt 2: đọc kỹ}

Danh sách đánh số:
\begin{enumerate}
  \item Đọc toàn bộ, bỏ qua chứng minh chi tiết.
  \item Ghi lại các thuật ngữ chưa hiểu.
  \begin{enumerate}
    \item Danh sách con tự đánh a, b, c.
  \end{enumerate}
\end{enumerate}

\begin{keyidea}
Sau lượt 2, cần tóm tắt được đóng góp chính của paper trong ba câu.
\end{keyidea}

\begin{question}[Câu hỏi khi đọc]
Paper giải quyết vấn đề gì? Baseline là gì? Kết quả có đáng tin không?
\end{question}

% =====================================================================
\section{Cách viết paper}

\subsection{Cấu trúc chuẩn}

Bảng đơn giản:

\begin{table}[H]
  \centering
  \caption{Các phần của một paper}
  \label{tab:structure}
  \begin{tabularx}{\linewidth}{l X}
    \toprule
    \textbf{Phần} & \textbf{Nội dung} \\
    \midrule
    Abstract     & Tóm tắt vấn đề, phương pháp, kết quả \\
    Introduction & Động lực, đóng góp chính \\
    Method       & Mô tả chi tiết phương pháp \\
    Experiments  & Thiết lập, kết quả, ablation \\
    Conclusion   & Tổng kết, hạn chế, hướng tiếp theo \\
    \bottomrule
  \end{tabularx}
\end{table}

Tham chiếu bảng: xem Bảng~\ref{tab:structure}.

\subsection{Công thức}

Công thức trong dòng: $\mathcal{L} = -\sum_i y_i \log \hat{y}_i$. Công thức riêng dòng:
\begin{equation}
  \theta_{t+1} = \theta_t - \eta \nabla_\theta \mathcal{L}(\theta_t)
  \label{eq:sgd}
\end{equation}

% =====================================================================
\section{Cách tái hiện paper}

\subsection{Danh sách kiểm tra}

\begin{checklist}
  \item Tìm code gốc và checkpoint (nếu có).
  \item Ghi lại phiên bản thư viện, GPU, seed.
  \item Tái hiện một bảng kết quả chính trước.
\end{checklist}

\begin{warning}
Kết quả lệch so với paper là bình thường. Cần ghi lại mức lệch và nguyên nhân dự đoán.
\end{warning}

\subsection{Chèn ảnh}

Chép ảnh vào thư mục \texttt{images/}, rồi gọi bằng tên file (không cần đuôi):

\begin{figure}[H]
  \centering
  \includegraphics[width=0.6\linewidth]{example-image}  % thay bằng tên ảnh, ví dụ: {loss-curve}
  \caption{Ví dụ chèn ảnh}
  \label{fig:example}
\end{figure}

Tham chiếu ảnh: xem Hình~\ref{fig:example}.

\subsection{Khung code}

Dán code trực tiếp:

\begin{lstlisting}[language=Python, caption={Vòng lặp huấn luyện}]
for epoch in range(num_epochs):
    for x, y in loader:
        loss = criterion(model(x), y)  # tính loss
        optimizer.zero_grad()
        loss.backward()
        optimizer.step()
\end{lstlisting}

Hoặc lấy code từ file trong thư mục \texttt{code/}:

\lstinputlisting[language=Python, caption={code/example.py}]{code/example.py}

Lệnh terminal:

\begin{lstlisting}[language=bash, numbers=none]
git clone https://github.com/user/repo.git
pip install -r requirements.txt
\end{lstlisting}

Code trong dòng: \lstinline[language=Python]|torch.manual_seed(42)|.

\todo{viết tiếp}

% =====================================================================
\bibliographystyle{plainnat}
\bibliography{references}

\end{document}
